\documentclass[10pt,a4paper]{article}

% ---------- lingua, codifica, font (testo + matematica) ----------
\usepackage[T1]{fontenc}
\usepackage[utf8]{inputenc}
\usepackage[italian]{babel}
\usepackage{mathpazo}                 % Palatino per testo e matematica
\usepackage{microtype}

% ---------- matematica ----------
\usepackage{amsmath,amssymb,amsfonts,mathtools,bm,cancel,mathrsfs,nicefrac}
\DeclareMathOperator{\sen}{sen}       % seno all'italiana
\DeclareMathOperator{\tg}{tg}
\DeclareMathOperator{\ctg}{cotg}
\DeclareMathOperator{\arcsen}{arcsen}
\DeclareMathOperator{\arctg}{arctg}
\newcommand{\dd}{\mathrm{d}}                       % differenziale
\newcommand{\vv}[1]{\vec{#1}}                      % vettore con freccia
\newcommand{\norma}[1]{\left\lvert#1\right\rvert} % modulo di vettore
\newcommand{\pedice}[1]{_{\text{#1}}}              % pedice testuale: E\pedice{in}

% ---------- unità di misura e numeri (virgola decimale) ----------
\usepackage{siunitx}
\sisetup{output-decimal-marker={,},per-mode=symbol,
         inter-unit-product=\ensuremath{{}\cdot{}},group-digits=false}

% ---------- grafici e figure ----------
\usepackage{tikz}
\usetikzlibrary{arrows.meta,calc,positioning,patterns,decorations.pathmorphing}
\usepackage{pgfplots}
\pgfplotsset{compat=1.16}
\usepackage{graphicx}
% Grafici collocati nel punto pertinente dello svolgimento.
\newenvironment{grafico}{\begin{figure}[h!]\centering}{\end{figure}}

% ---------- impaginazione ----------
\usepackage[a4paper,left=1.8cm,right=1.8cm,top=2.3cm,bottom=2.2cm]{geometry}
\usepackage{fancyhdr,lastpage,needspace,enumitem,booktabs,array,colortbl,tabularx}
\setlength{\parindent}{0pt}
\setlength{\parskip}{3pt}
\setlength{\headheight}{14pt}
\setlist[enumerate]{leftmargin=1.7em,itemsep=2pt,topsep=2pt,parsep=0pt}
\setlist[itemize]{leftmargin=1.3em,itemsep=2pt,topsep=2pt,parsep=0pt}

% ---------- colori degli indicatori ----------
\usepackage[table,dvipsnames]{xcolor}
\definecolor{anac}{HTML}{C62828}\definecolor{anabg}{HTML}{FDECEA} % rosso
\definecolor{svic}{HTML}{1A56DB}\definecolor{svibg}{HTML}{E8F0FE} % blu
\definecolor{argc}{HTML}{A67C00}\definecolor{argbg}{HTML}{FFF6CC} % giallo
\definecolor{comc}{HTML}{2E7D32}\definecolor{combg}{HTML}{E6F4EA} % verde
% Solo griglie d'esame (su richiesta): IND arancio, DAT teal
\definecolor{indc}{HTML}{E65100}\definecolor{indbg}{HTML}{FFF0E0}
\definecolor{datc}{HTML}{00796B}\definecolor{datbg}{HTML}{E0F2F1}

% Riferimenti agli indicatori: usare queste macro in tutto il documento.
\newcommand{\coloreANA}{anac}\newcommand{\nomeANA}{Analizzare}
\newcommand{\coloreSVI}{svic}\newcommand{\nomeSVI}{Sviluppare e calcolare}
\newcommand{\coloreARG}{argc}\newcommand{\nomeARG}{Argomentare}
\newcommand{\coloreCOM}{comc}\newcommand{\nomeCOM}{Comunicare}
\newcommand{\coloreIND}{indc}\newcommand{\nomeIND}{Individuare}
\newcommand{\coloreDAT}{datc}\newcommand{\nomeDAT}{Interpretare i dati}
% Griglie d'esame: matematica \renewcommand{\nomeANA}{Comprendere}, \nomeSVI{Sviluppare il processo risolutivo};
% fisica/mista: \nomeSVI{Sviluppare il processo risolutivo}, \nomeDAT{Interpretare criticamente i dati} (mista: ...rappresentare, elaborare).
\newcommand{\sigla}[1]{\textcolor{\csname colore#1\endcsname}{\textbf{#1}}}
\newcommand{\indicatore}[1]{\textcolor{\csname colore#1\endcsname}{\textbf{#1: \csname nome#1\endcsname}}}

% ---------- box colorati e chip ----------
\usepackage{tcolorbox}
\tcbuselibrary{skins,breakable}

% Box passaggio chiave: \begin{boxana}[Titolo] ... \end{boxana}  (ana|svi|arg|com)
\newcommand{\defbox}[4]{%
  \newtcolorbox{box#1}[1][Passaggio chiave]{enhanced,breakable,frame hidden,
    borderline west={2.5pt}{0pt}{#2},colback=#3,boxrule=0pt,arc=1.5pt,
    left=7pt,right=5pt,top=2pt,bottom=2pt,before skip=4pt,after skip=4pt,
    fontupper=\small,before upper={\textbf{\idevidenza}\enspace\sigla{#4}\enspace\textbf{##1}\par}}}
\defbox{ana}{anac}{anabg}{ANA}
\defbox{svi}{svic}{svibg}{SVI}
\defbox{arg}{argc}{argbg}{ARG}
\defbox{com}{comc}{combg}{COM}
\defbox{ind}{indc}{indbg}{IND}
\defbox{dat}{datc}{datbg}{DAT}

% Box enunciato
\newtcolorbox{enun}{colback=black!3,colframe=black!25,boxrule=0.4pt,arc=2pt,
  left=6pt,right=6pt,top=3pt,bottom=3pt,fontupper=\small,
  before skip=5pt,after skip=5pt}

% Chip identificativo e massimo in punti; i dati sono forniti dal registro approvato.
\newcommand{\chip}[4]{\tcbox[on line,boxsep=0pt,left=3pt,right=3pt,top=1pt,bottom=1pt,
  arc=4pt,boxrule=0.4pt,colframe=#1,colback=#2,fontupper=\scriptsize\bfseries\color{#1}]{#3\ #4}}
\newcommand{\ANA}[2]{\chip{anac}{anabg}{\sigla{ANA}}{#1\, / massimo #2\,p}}
\newcommand{\SVI}[2]{\chip{svic}{svibg}{\sigla{SVI}}{#1\, / massimo #2\,p}}
\newcommand{\ARG}[2]{\chip{argc}{argbg}{\sigla{ARG}}{#1\, / massimo #2\,p}}
\newcommand{\COM}[2]{\chip{comc}{combg}{\sigla{COM}}{#1\, / massimo #2\,p}}
\newcommand{\IND}[2]{\chip{indc}{indbg}{\sigla{IND}}{#1\, / massimo #2\,p}}
\newcommand{\DAT}[2]{\chip{datc}{datbg}{\sigla{DAT}}{#1\, / massimo #2\,p}}
\newcommand{\ess}{\ensuremath{\bigstar}}
\newcommand{\legenda}{\indicatore{ANA}\quad\indicatore{SVI}\par
  \indicatore{ARG}\quad\indicatore{COM}}
% Griglie d'esame: \legendaesame{ANA}{IND}{SVI}{ARG} (matematica) o \legendaesame{ANA}{SVI}{DAT}{ARG} (fisica, mista)
\newcommand{\legendaesame}[4]{\indicatore{#1}\quad\indicatore{#2}\par\indicatore{#3}\quad\indicatore{#4}}
% Limite categorico sull'intera verifica; il registro non incrementa il conteggio.
\newcounter{numeroevidenze}
\newcommand{\controllaevidenze}{\stepcounter{numeroevidenze}%
  \ifnum\value{numeroevidenze}>40\relax
    \PackageError{correttore}{Superato il limite categorico di 40 evidenze}{Accorpare le evidenze prima di generare il correttore.}%
  \fi}
% Metadati mantenuti nel sorgente; nel corpo il titolo viene emesso solo dal box.
\newcommand{\idevidenza}{}
\newcommand{\evidenza}[5]{\controllaevidenze\def\idevidenza{#1}\needspace{8\baselineskip}}
% Registro raggruppato: esercizio, argomento comune, righe ID/nome/punto/indicatore/massimo.
\newcommand{\registrogruppo}[3]{\needspace{6\baselineskip}\par\smallskip{\small
  \textbf{Esercizio #1}\enspace #2\par
  \begin{tabularx}{\linewidth}{@{}lXllr@{}}\toprule
  ID & Evidenza & Punto & Ind. & Max\\\midrule
  #3\bottomrule\end{tabularx}\par}}
\newcommand{\inizioesercizi}{\clearpage}
\newcommand{\riepiloghi}[2]{\sez{Riepiloghi}
  \noindent\begin{minipage}[t]{0.58\linewidth}#1\end{minipage}\hfill
  \begin{minipage}[t]{0.39\linewidth}#2\end{minipage}\par}

\newcommand{\livelli}[5]{\begin{description}[leftmargin=2.7cm,style=nextline,font=\small\bfseries]
\item[Non raggiunto] #1
\if\relax\detokenize{#2}\relax\else\item[Iniziale] #2\fi
\item[Base] #3
\if\relax\detokenize{#4}\relax\else\item[Intermedio] #4\fi
\item[Avanzato] #5
\end{description}}

% Intestazioni di colonna colorate
\newcommand{\hANA}{\cellcolor{anabg}\sigla{ANA}}
\newcommand{\hSVI}{\cellcolor{svibg}\sigla{SVI}}
\newcommand{\hARG}{\cellcolor{argbg}\sigla{ARG}}
\newcommand{\hCOM}{\cellcolor{combg}\sigla{COM}}
\newcommand{\hIND}{\cellcolor{indbg}\sigla{IND}}
\newcommand{\hDAT}{\cellcolor{datbg}\sigla{DAT}}

% ---------- macro di struttura ----------
% Sezione semplice
\newcommand{\sez}[1]{\needspace{4\baselineskip}\par\vspace{1em}%
  {\large\bfseries #1}\par\vspace{-0.4em}\rule{\linewidth}{0.6pt}\par\vspace{0.2em}}
% Esercizio: \esercizio{Titolo}{ANA}{SVI}{ARG}{COM}{TOT}
\newcommand{\esercizio}[6]{\needspace{8\baselineskip}\par\vspace{1em}%
  {\large\bfseries #1}\par\vspace{0.25em}}
% Punto: \punto{a) titolo}
\newcommand{\punto}[1]{\needspace{6\baselineskip}\par\vspace{0.7em}{\bfseries #1}\par\vspace{0.2em}}
% Tabella punteggi del punto: \tabpunti{ANA}{SVI}{ARG}{COM}{TOT}  (usare 0 se assente)
\newcommand{\tabpunti}[5]{\par\noindent\renewcommand{\arraystretch}{1.15}\small
  \begin{tabular}{*{5}{|>{\centering\arraybackslash}p{1.4cm}}|}\hline
  \hANA & \hSVI & \hARG & \hCOM & \textbf{Tot}\\\hline
  #1 & #2 & #3 & #4 & \textbf{#5}\\\hline\end{tabular}\par\smallskip}
% Griglie d'esame: \tabpuntiesame{ANA}{IND}{SVI}{ARG}{v1}{v2}{v3}{v4}{TOT} (sigle nell'ordine dei valori)
\newcommand{\tabpuntiesame}[9]{\par\noindent\renewcommand{\arraystretch}{1.15}\small
  \begin{tabular}{*{5}{|>{\centering\arraybackslash}p{1.4cm}}|}\hline
  \csname h#1\endcsname & \csname h#2\endcsname & \csname h#3\endcsname & \csname h#4\endcsname & \textbf{Tot}\\\hline
  #5 & #6 & #7 & #8 & \textbf{#9}\\\hline\end{tabular}\par\smallskip}
% Nota (alternative, errori trascinati, criteri)
\newcommand{\nota}[1]{\begin{itemize}[font=\footnotesize]\item {\footnotesize #1}\end{itemize}}
\newcommand{\criteri}[1]{\par\nopagebreak[3]\smallskip\needspace{3\baselineskip}%
  {\footnotesize\textbf{Alternative, errori e criteri}\par
  \begin{itemize}[itemsep=1pt,topsep=1pt]#1\end{itemize}}}

% ---------- intestazione e piè di pagina ----------
\newcommand{\materia}{}
\newcommand{\classe}{}
\newcommand{\argomenti}{}
\pagestyle{fancy}\fancyhf{}
\renewcommand{\headrulewidth}{0.3pt}
\lhead{\small\textsc{Correttore\ \materia}}
\rhead{}
\cfoot{\small Pag.\ \thepage\ di \pageref{LastPage}}

\renewcommand{\materia}{Matematica}
\renewcommand{\classe}{Classe quinta, liceo scientifico \textperiodcentered\ Prova modellata sulla seconda prova di matematica}
\renewcommand{\nomeANA}{Comprendere}\renewcommand{\nomeSVI}{Sviluppare il processo risolutivo}
\begin{document}
\begin{center}
{\LARGE\bfseries Correttore\ \materia}\\[3pt]
\classe
\end{center}

\legendaesame{ANA}{IND}{SVI}{ARG}
\sez{Registro delle evidenze}
\registrogruppo{1}{FUN, INT}{
E01 & Individua dominio e limiti da calcolare & intero & \sigla{ANA} & 50\\
E02 & Calcola limiti e asintoto & intero & \sigla{SVI} & 50\\
E03 & Sceglie il segno della derivata prima & intero & \sigla{IND} & 80\\
E04 & Calcola $f'$ e individua il massimo & intero & \sigla{SVI} & 80\\
E05 & Sceglie la derivata seconda & intero & \sigla{IND} & 70\\
E06 & Calcola $f''$ e il flesso & intero & \sigla{SVI} & 50\\
E07 & Traduce l'area in un integrale & intero & \sigla{ANA} & 100\\
E08 & Sceglie l'integrazione per parti & intero & \sigla{IND} & 70\\
E09 & Calcola l'integrale & intero & \sigla{SVI} & 70\\
E10 & Commenta la coerenza con il grafico & intero & \sigla{ARG} & 100\\
}
\registrogruppo{2}{TEO}{
E11 & Individua le proprietà utili & intero & \sigla{ANA} & 100\\
E12 & Sceglie valori intermedi e monotonia & intero & \sigla{IND} & 80\\
E13 & Motiva e localizza le due soluzioni & intero & \sigla{ARG} & 100\\
}
\riepiloghi{%
\begin{tabular}{@{}lrrrrr@{}}\toprule
 & \hANA & \hIND & \hSVI & \hARG & Tot\\\midrule
Esercizio 1 & 150 & 220 & 250 & 100 & 720\\
Esercizio 2 & 100 & 80 & 0 & 100 & 280\\
\midrule
Totale & 250 & 300 & 250 & 200 & 1000\\\bottomrule
\end{tabular}}{%
\begin{tabular}{@{}lr@{}}\toprule
Argomento & Max\\\midrule
FUN, Studio di funzione & 480\\
INT, Integrali & 240\\
TEO, Teoremi sulle funzioni continue & 280\\
\midrule
Totale & 1000\\\bottomrule
\end{tabular}}

\inizioesercizi
\esercizio{Esercizio 1 -- Problema}{}{}{}{}{}
\begin{enun}
Sia $f(x)=x\,e^{-x}$. (a) Determina dominio, limiti agli estremi e asintoti. (b) Studia monotonia ed estremi. (c) Studia concavità e flessi. (d) Calcola l'area della regione fra il grafico, l'asse $x$ e le rette $x=0$ e $x=1$. (e) Commenta la coerenza dei risultati con l'andamento del grafico.
\end{enun}
\tabpuntiesame{ANA}{IND}{SVI}{ARG}{150}{220}{250}{100}{720}
\evidenza{E01}{1}{ANA}{FUN}{50}
\begin{boxana}[Individua dominio e limiti da calcolare]
\begin{enumerate}
\item La funzione è definita e continua su $\mathbb R$: nessuna condizione di esistenza.
\item I limiti rilevanti sono per $x\to-\infty$ e $x\to+\infty$.
\end{enumerate}
\end{boxana}
\evidenza{E02}{1}{SVI}{FUN}{50}
\begin{boxsvi}[Calcola limiti e asintoto]
\begin{enumerate}
\item $\lim_{x\to-\infty}xe^{-x}=-\infty$.
\item $\lim_{x\to+\infty}xe^{-x}=0$ (l'esponenziale prevale): asintoto orizzontale $y=0$.
\end{enumerate}
\end{boxsvi}
\evidenza{E03}{1}{IND}{FUN}{80}
\begin{boxind}[Sceglie il segno della derivata prima]
\begin{enumerate}
\item Monotonia ed estremi si studiano con il segno di $f'$.
\item I punti stazionari si trovano con $f'(x)=0$.
\end{enumerate}
\end{boxind}
\evidenza{E04}{1}{SVI}{FUN}{80}
\begin{boxsvi}[Calcola $f'$ e individua il massimo]
\begin{enumerate}
\item $f'(x)=(1-x)e^{-x}$.
\item $f'>0$ per $x<1$ e $f'<0$ per $x>1$.
\item Massimo assoluto in $\left(1,\frac1e\right)$.
\end{enumerate}
\end{boxsvi}
\evidenza{E05}{1}{IND}{FUN}{70}
\begin{boxind}[Sceglie la derivata seconda]
\begin{enumerate}
\item Concavità e flessi si studiano con il segno di $f''$.
\end{enumerate}
\end{boxind}
\evidenza{E06}{1}{SVI}{FUN}{50}
\begin{boxsvi}[Calcola $f''$ e il flesso]
\begin{enumerate}
\item $f''(x)=(x-2)e^{-x}$.
\item Flesso in $\left(2,\frac2{e^2}\right)$: concavità verso il basso per $x<2$, verso l'alto per $x>2$.
\end{enumerate}
\end{boxsvi}
\evidenza{E07}{1}{ANA}{INT}{100}
\begin{boxana}[Traduce l'area in un integrale]
\begin{enumerate}
\item Su $[0,1]$ è $f\ge0$: l'area è $\int_0^1 xe^{-x}\,dx$.
\end{enumerate}
\end{boxana}
\evidenza{E08}{1}{IND}{INT}{70}
\begin{boxind}[Sceglie l'integrazione per parti]
\begin{enumerate}
\item Il prodotto $x\cdot e^{-x}$ si integra per parti con $u=x$ e $dv=e^{-x}dx$.
\end{enumerate}
\end{boxind}
\evidenza{E09}{1}{SVI}{INT}{70}
\begin{boxsvi}[Calcola l'integrale]
\begin{enumerate}
\item $\int xe^{-x}dx=-xe^{-x}-e^{-x}+c$.
\item $\int_0^1xe^{-x}dx=1-\frac2e\approx0{,}264$.
\end{enumerate}
\end{boxsvi}
\evidenza{E10}{1}{ARG}{FUN}{100}
\begin{boxarg}[Commenta la coerenza con il grafico]
\begin{enumerate}
\item Il massimo vale $\frac1e\approx0{,}37$: su $[0,1]$ l'area è minore di $0{,}37$, coerente con $0{,}264$.
\item Crescita fino a $x=1$, flesso in $x=2$ e tendenza a $0$ concordano con i segni di $f'$, $f''$ e con l'asintoto.
\end{enumerate}
\end{boxarg}
\esercizio{Esercizio 2 -- Quesito}{}{}{}{}{}
\begin{enun}
Dimostra che l'equazione $x\,e^{-x}=\tfrac14$ ha esattamente due soluzioni reali e localizzale in intervalli di ampiezza $1$.
\end{enun}
\tabpuntiesame{ANA}{IND}{SVI}{ARG}{100}{80}{0}{100}{280}
\evidenza{E11}{2}{ANA}{TEO}{100}
\begin{boxana}[Individua le proprietà utili]
\begin{enumerate}
\item $f$ cresce su $(-\infty,1]$ da $-\infty$ a $\frac1e$ e decresce su $[1,+\infty)$ da $\frac1e$ a $0$.
\item Per $x\le0$ è $f\le0<\frac14$: le soluzioni sono in $x>0$.
\end{enumerate}
\end{boxana}
\evidenza{E12}{2}{IND}{TEO}{80}
\begin{boxind}[Sceglie valori intermedi e monotonia]
\begin{enumerate}
\item Su ogni tratto monotono la continuità assicura l'esistenza e la monotonia l'unicità.
\end{enumerate}
\end{boxind}
\evidenza{E13}{2}{ARG}{TEO}{100}
\begin{boxarg}[Motiva e localizza le due soluzioni]
\begin{enumerate}
\item $f(0)=0<\frac14<f(1)=\frac1e$: una soluzione in $(0,1)$, unica per la crescita.
\item $f(2)=\frac2{e^2}\approx0{,}271>\frac14>f(3)=\frac3{e^3}\approx0{,}149$: una soluzione in $(2,3)$, unica per la decrescita.
\item Non ce ne sono altre.
\end{enumerate}
\end{boxarg}
\end{document}
